\section{Engine limits and the solar floor}
\label{sec:constraints}
\subsection{The bounded-engine selector}
For normal running cost $\ell_0+\ell_2|u|^2$, set
$C=p_b+\ell_2$. Fixed-time effort uses $(\ell_0,\ell_2)=(0,1)$;
minimum time uses $(1,0)$. The Hamiltonian is
\begin{equation}
 H=p_r\cdot v+p_v\cdot g-\ell_0+p_v\cdot u-C|u|^2.
\end{equation}
Let $s=|p_v|$, $a_-=w/\nu_+$ and $a_+=w/\nu_-$. For $C>0$,
\begin{equation}
 \widehat a=\operatorname{clip}(s/(2C),a_-,a_+),\qquad
 \mathcal G=\max(0,s\widehat a-C\widehat a^2).
 \label{eq:selector}
\end{equation}
Burn in direction $p_v/s$ when the powered gain is positive and coast
when it is negative. At a tie both maximize. For $C\le0$, compare
coast with $a_+$. If $s=0$ and a burn wins, its direction is not
determined by this pointwise condition.

For $C>0$ the four regular magnitude regions are
\begin{equation}
 a^*=\begin{cases}
 0&s<Ca_-,\\
 a_-&Ca_-<s<2Ca_-,\\
 s/(2C)&2Ca_-<s<2Ca_+,\\
 a_+&s>2Ca_+.
 \end{cases}
\end{equation}
The lowest powered thrust uses maximum exhaust; the highest uses
minimum exhaust. State and costate evolution determines the order of
these modes along a transfer.

The kinematic adjoints retain \eqref{eq:idealPMP}, but the resource
adjoint changes on a fixed-exhaust boundary:
\begin{equation}
 \begin{gathered}
 C'=\frac{s-2Ca}{\nu}\quad(a=w/\nu),\\
 C'=0\quad\hbox{on interior burns and coasts}.
 \end{gathered}
 \label{eq:Cprime}
\end{equation}
Indeed $\partial a/\partial b=-1/\nu$, so differentiation of the
optimized Hamiltonian gives \eqref{eq:Cprime}. Ordinary switches
inside $b>0$ have continuous states and costates. Dry-mass exhaustion
is a state-constraint event and requires its additional multiplier.
For free inactive terminal reserve in an isolated effort problem,
$p_b(h)=0$; an embedded leg receives its resource covector from the
next leg.

Autonomy and rotational symmetry preserve $H^*$ and $\mathcal L$
on the regular phase. The product $\lambda_m m^2$ does not stay
constant on capped arcs, because the acceleration bounds depend on
mass. Appendix~\ref{app:repair} gives the one-dimensional propagation. There is no derived universal finite
switching sequence for the Kepler problem.

With ordinary ballistic endpoint ephemerides, the bounded counterparts
of \eqref{eq:dates} are
\begin{equation}
 F_h=-\mathcal G_f,\qquad
 F_d=\mathcal G_0-\mathcal G_f,
 \label{eq:boundeddate}
\end{equation}
when terminal resource variations vanish or are included through their
transversality conditions. The ideal case has $p_b=0$ and
$\mathcal G=|u|^2$.

\subsection{Auditing an ideal trajectory}
Given an endpoint-exact ideal path, write
$I(t)=\int_0^t|u|^2ds$ and $w_I=\chi+I$. Its actual full-power exhaust
on a burn is $\nu_e=w_I/|u|$. It is band-feasible precisely when
$I(h)\le\kappa$ and
\begin{equation}
 w_I/\nu_+\le|u|\le w_I/\nu_-
 \label{eq:bandaudit7}
\end{equation}
almost everywhere on powered intervals; exact coasts pass separately.
Certified interval bounds can establish \eqref{eq:bandaudit7}
between samples. Clipping a violating control changes its endpoint.
If a globally ideal-minimal path passes this audit, it is also globally
minimal in the restricted problem at that time.

Continuous throttling is a different engine model. For duty fraction
$\delta$, $a=\delta w/\nu$ and $w'=a^2/\delta$.
The fuel-minimal realization of a prescribed acceleration history uses
\begin{equation}
 \nu=\min(\nu_+,w/a),\qquad
 w'=\max(a^2,wa/\nu_+).
 \label{eq:throttle}
\end{equation}
It also requires $a\le w/\nu_-$ and $w(h)\le\chi_d$.
This scalar equation preserves the prescribed spatial path exactly
under continuous throttling. An on/off approximation must preserve
both mean acceleration and mass flow: a duty cycle has mean effort
rate $\delta a_{\rm burn}^2$, not $(\delta a_{\rm burn})^2$.
Exact endpoint correction for finite pulses is an additional problem.

\subsection{A hard radial constraint}
Set $R_c=r_{\min}/R_0$ and impose
$c=(R_c^2-|r|^2)/2\le0$. For the solar model
$r_{\min}=0.1$ AU; $R_c=0.1$ only when $R_0=1$ AU.
With radial and tangential speeds $y,z$,
\begin{equation}
 \begin{gathered}
 R'=y,\qquad
 y'=z^2/R-1/R^2+u_r,\\
 \theta'=z/R,\qquad z'=u_\theta-yz/R.
 \end{gathered}
\end{equation}
The constraint has relative order two:
$c'=-r\cdot v$ and
$c''=-|v|^2+1/R-r\cdot u$.
At a boundary interval,
\begin{equation}
 \begin{gathered}
 R=R_c,\quad y=0,\quad z'=u_\theta,\\
 u_r=A(z):=R_c^{-2}-z^2/R_c.
 \end{gathered}
 \label{eq:boundary}
\end{equation}
A circular coast with $z=\pm R_c^{-1/2}$ is a zero-effort globally
minimal boundary arc for its own endpoints. Thus boundary intervals
must be permitted.

Using the original measure-costate convention, a nonnegative measure
$dM$ is supported on the contact set \cite{HartlSethiVickson1995,BonnansHermant2009}:
\begin{equation}
 dp_r=-Dg^Tp_v\,dt-r\,dM,\qquad dp_v=-p_r\,dt.
\end{equation}
At an atom $\gamma\ge0$,
$p_r^+=p_r^--\gamma r$ and $p_v^+=p_v^-$.
The ideal thrust is continuous, while its derivative can jump.
Tangency makes the Hamiltonian continuous; rotational symmetry
preserves $\mathcal L$. Higher-order transformed adjoints use a
different representation and should not be mixed with these jumps.

For the ideal boundary problem, substituting \eqref{eq:boundary}
gives $J_B=\int(A(z)^2+z'^2)dt$ and
\begin{equation}
 z''=-\frac{\mathcal L}{2R_c}-\frac{2A(z)z}{R_c}.
\end{equation}
The original radial adjoint fixes its measure density:
\begin{equation}
 \mu_M=
 \frac{-4A/R_c^3-6z'^2/R_c-8zz''/R_c-2Az^2/R_c^2}{R_c}
 \ge0.
 \label{eq:density}
\end{equation}
For completeness, on the circle
$p_v=2Ae_r+2z'e_\theta$ and $e_r'=(z/R_c)e_\theta$ give
$p_r=-p_v'=6zz'e_r/R_c-2(z''+Az/R_c)e_\theta$.
The radial component of $p_r'=-Dg^Tp_v-R_ce_r\mu_M$, with
$Dg\,e_r=2e_r/R_c^3$, gives \eqref{eq:density} term by term.
This sign test and the entry/exit conditions are part of the solution.
For constant $z$, a subcircular boundary burn has $\mu_M<0$ and
fails the normal necessary conditions.

With engine bounds the boundary is feasible only if $|A|\le a_+$ and
\begin{equation}
 \sqrt{\max(0,a_-^2-A^2)}
 \le |u_\theta|\le\sqrt{a_+^2-A^2}.
\end{equation}
Coast is separately possible when $A=0$. Add the radial measure to
the bounded-engine adjoints and retain \eqref{eq:Cprime}.

A smooth radial penalty preserves time and rotation symmetry but
changes the objective and its sensitivities. It does not in general
enforce the hard floor. The barrier $\epsilon/(R-R_c)^2$ on
$R>R_c$ excludes contact at finite penalized cost and can initialize
continuation, but cannot itself represent a boundary interval.
For a regular constrained minimum with fixed endpoints,
\begin{equation}
 \frac{\partial J^*}{\partial R_c}=R_cM([0,h])\ge0.
\end{equation}
The physical effort must be reported separately from any penalty cost.

For runtime safety, a cubic Bezier position interpolant $p$ with
controls $B_i$ has squared-radius Bernstein coefficients
\begin{equation}
 c_k=\sum_{i+j=k}
 \frac{\binom3i\binom3j}{\binom6k}B_i\cdot B_j.
\end{equation}
If a validated trajectory error is $|r-p|\le E$, then
$\min_k c_k\ge(R_c+E)^2$ is a sufficient segment pass.
Subdivide unresolved segments. The error enclosure must include
trajectory and arithmetic error, not just interpolation of sampled data.
